\section{简易程序与速裁程序}
\begin{enumerate}
  \item 适用范围。\bs{清楚 + 认罪 + 被告人同意}。
        \begin{enumerate}
          \item 认罪。对指控的\zd{犯罪事实}无异议。区分对行为性质的辩解。
                \al{朱某承认检察院指控的犯罪事实，但认为自己应当构成抢夺罪，不该构成抢劫罪 \dui}。
          \item 不认罪。扩大解释。可能无罪、无罪辩护、\bs{共犯部分不认罪}。
          \item 例外。\zd{盲聋哑，半疯傻，影响大}。
        \end{enumerate}
  \item 程序。
        \begin{enumerate}
          \item 申请 \hz 检察院建议 \yjt 法院决定。
          \item \zd{基层一审、3 年为界}、\bs{一般应当}当庭宣判（1 次搞不定）。
          \item 3 年以下可合可独，3 年以上应当合议庭。
          \item 3 年以下 20 日，3 年以上一个半月。
          \item 检察院应当派员出席。有辩护人的，应当通知出庭。
          \item 庭审不能简化。\bs{必须当庭询问被告人对指控的犯罪事实的意见}。
          \item 转普通。从适用条件推。没有普通转简易。
        \end{enumerate}
  \item \bs{公诉}速裁程序。比简易更简易。\bs{+ 轻微（3 年以下）+ 附民已解决}。
        \begin{enumerate}
          \item 一人独任、10 + 5、\bs{应当}当庭宣判。
          \item \zd{庭前准备时间不受限制}。
          \item 转普通。从适用条件推。
        \end{enumerate}
\end{enumerate}
